\newcommand{\ModuleID}{EL-II.10}
\newcommand{\ModuleTitle}{Relative Clauses of Characteristic and Related Uses}
\newcommand{\ModuleStage}{Stage II — Core Grammar}
\newcommand{\ModuleUnit}{II.10}
\newcommand{\ModuleOutcome}{Distinguish factual, characteristic, purpose, causal, and concessive relatives through an evidence scale while parsing the relative form independently}
\newcommand{\ModulePrerequisites}{EL-II.9 mastery: classify and transform conditional stance and time}
\newcommand{\ModuleExtensionPractice}{Worksheets II.38 and II.39 remain optional remediation or delayed-recall variants}
\newcommand{\ModuleNext}{EL-II.11, Commands, Prohibitions, Wishes, and Fearing}
\newcommand{\ModuleMastery}{Choose and defend mood in at least nine of ten contrasting relatives, parse every relative word's case correctly, and compose four clause types whose English renderings preserve their intended differences.}
\newcommand{\ModuleDelayNote}{}
\newcommand{\ModuleLessonSource}{\CourseRoot/shared/content/volume2/all}
\newcommand{\ModuleExerciseSource}{\CourseRoot/shared/exercises/volume2/all}
\newcommand{\ModuleWorkedBank}{II}
\newcommand{\ModuleExerciseBank}{II}
\newcommand{\ModuleWorksheetVariants}{A,D}
\newcommand{\ModuleScopeNote}{This packet uses an evidence scale rather than claiming that every subjunctive relative has one English formula.}
\newcommand{\ModulePractice}{%
  \SubmoduleExtension
    {Fact or characteristic?}
    {An evidence scale makes the factual/characteristic decision cumulative rather than mechanical.}
    {Pair lesson relatives that differ in mood or antecedent type. Score evidence from identified versus indefinite/class antecedent, mood, negation, governing expression, and discourse purpose; then state the narrowest defensible reading.}
    {Give two independent signals that support a characteristic reading.}
    {A negative, indefinite, or class-defining antecedent and subjunctive mood commonly converge on characteristic; discourse that describes a kind rather than reports an individual fact strengthens it. No one signal excuses parsing the relative pronoun and verb.}
  \SubmoduleExtension
    {Reading the forms in context}
    {Relative purpose, cause, and concession answer different discourse questions even when subjunctive mood recurs.}
    {For each related-use example in the lesson, write the question it answers—what end, what alleged reason, despite what circumstance—and cite the governor or context that supports the label.}
    {Explain why characteristic is not the default label for every subjunctive relative.}
    {Each use requires its own relation and evidence. Characteristic defines a class or kind; purpose expresses intended end; cause supplies or attributes reason; concession presents an obstacle that does not prevent the main assertion.}
  \SubmoduleExtension
    {Writing laboratory}
    {Minimal pairs reveal how mood and viewpoint change while relative morphology can remain stable.}
    {Complete the lesson's indicative/subjunctive identity pair and the requested purpose and characteristic clauses. Label antecedent coordinates, internal case, mood, and discourse relation on each.}
    {State which relative-pronoun properties can stay unchanged when clause use changes.}
    {With the same antecedent and internal function, gender, number, and case can remain unchanged; mood and discourse relation change. Each composed clause must preserve the intended use and a faithful English distinction.}
}
